\chapter{People and Conduct Risk Assessment}
\label{chap:People and Conduct Risk Assessment}

%---------------------------- SECTION ----------------------------
\section{The Risk That Wears a Badge}
\paragraph{Every risk covered in this book — every cyber attack that succeeds, every regulatory breach that occurs, every operational failure that disrupts, every reputational crisis that unfolds — has a human dimension. Behind every system vulnerability that goes unpatched is a person who decided it was not a priority. Behind every misselling scandal is a sales culture that rewarded volume over suitability. Behind every fraud is an individual who made a choice. Behind every whistleblowing failure is an organisation whose people did not feel safe speaking up.}
\paragraph{People are not merely the managers of risk — they are, in a profound and inescapable sense, its primary source. The behaviour, the judgement, the values, the competence, and the motivations of the people within an organisation determine, more than any other single factor, whether its risk management framework actually works or merely exists on paper.}
\paragraph{This insight is not new. The field of human factors — the systematic study of how human characteristics affect performance in complex systems — has been developing for decades, primarily in aviation, healthcare, and nuclear safety, where the consequences of human error are most immediately and most visibly catastrophic. But its application to compliance and legal risk management has been slower to develop — and the gap between the sophistication of the technical risk frameworks that organisations deploy and the attention they give to the human factors that determine whether those frameworks actually work remains wide in many organisations.}
\paragraph{\textbf{Conduct risk} — the risk that the behaviour of the organisation's people will cause harm to customers, markets, or the organisation itself — has emerged as a specific and prominent regulatory concern, particularly in financial services, following the wave of conduct failures that characterised the years before and after the 2008 financial crisis. The mis-selling of payment protection insurance, the manipulation of LIBOR and foreign exchange benchmarks, the widespread misselling of interest rate hedging products to small businesses — these were not primarily failures of systems or processes. They were failures of human behaviour, enabled and in many cases incentivised by cultures and structures that prioritised financial performance over customer welfare and market integrity.}
\paragraph{For compliance and legal professionals, people and conduct risk is therefore both a specific regulatory concern — directly addressed by frameworks such as the Senior Managers and Certification Regime — and a fundamental dimension of every other category of risk assessed in this book. This chapter provides the framework for understanding, assessing, and managing it.}

%---------------------------- SECTION ----------------------------
\section{Defining the Territory — People Risk and Conduct Risk}
\paragraph{Before exploring the assessment framework, it is worth drawing a clear distinction between two related but distinct concepts that are sometimes conflated.}
\paragraph{\textbf{People risk} is the broader category — encompassing all the ways in which the characteristics, behaviour, and circumstances of individuals within the organisation create risk. It includes:}
\begin{itemize}
    \item{The risk of human error — mistakes, oversights, and failures of judgement that occur in the normal course of work.}
    \item{The risk of inadequate competence — individuals who lack the knowledge, skills, or experience to perform their roles to the required standard.}
    \item{The risk of key person dependency — the organisation's over-reliance on specific individuals whose absence or departure creates significant operational or knowledge risk.}
    \item{The risk of workforce capacity — insufficient numbers of adequately skilled people to perform the organisation's activities safely and to the required standard.}
    \item{The risk of cultural failure — the organisation's values, norms, and incentive structures creating conditions in which poor behaviour is enabled or encouraged.}
\end{itemize}
\paragraph{\textbf{Conduct risk} is a specific subset of people risk — focused on behaviour that is, or has the potential to be, harmful. In financial services, conduct risk is typically defined as the risk that the behaviour of the firm or its staff leads to poor outcomes for clients, counterparties, or markets, or that it results in regulatory breach. More broadly, conduct risk encompasses:}
\begin{itemize}
    \bitem{Customer conduct risk}{behaviour that causes harm to customers — through misselling, poor advice, unfair treatment, or failure to act in customers' interests}
    \bitem{Market conduct risk}{behaviour that undermines the integrity of financial markets — through insider trading, market manipulation, benchmark manipulation, or other forms of market abuse.}
    \bitem{Internal conduct risk}{behaviour that harms the organisation itself or its stakeholders — through fraud, theft, unauthorised trading, or conflicts of interest that are not properly managed.}
    \bitem{Regulatory conduct risk}{behaviour towards regulators that is dishonest, obstructive, or otherwise inappropriate — including misleading regulatory disclosures, obstruction of regulatory investigations, or failure to comply with regulatory requirements.}
\end{itemize}
\paragraph{The distinction between people risk and conduct risk matters for assessment purposes — because they require somewhat different analytical approaches and different control responses. People risk is largely about capability, capacity, and process design. Conduct risk is primarily about behaviour, culture, and incentives.}

%---------------------------- SECTION ----------------------------
\section{The Regulatory Framework — Senior Managers and Certification Regime}
\paragraph{In financial services, the regulatory framework most directly relevant to people and conduct risk is the \textbf{Senior Managers and Certification Regime (SM\&CR)} — introduced by the FCA and PRA following the recommendations of the Parliamentary Commission on Banking Standards, and applying to banks, insurers, and most other FCA-regulated firms.}
\paragraph{The SM\&CR is built on a simple but powerful premise: that individual accountability — knowing clearly who is responsible for what, and that individuals will be held personally to account for the decisions they make and the cultures they create — is the most effective mechanism for driving improved conduct in financial services.}
\paragraph{The regime has three main components:}

\subsection{The Senior Managers Regime}
\paragraph{The \textbf{Senior Managers Regime} applies to the most senior individuals in regulated firms — those performing \textbf{Senior Management Functions (SMFs)} including the Chief Executive, Chief Finance Officer, Chief Risk Officer, Chief Compliance Officer, and other board and senior management roles. Senior Managers are required to:}
\begin{itemize}
    \item{Be approved by the FCA and/or PRA before taking up their role.}
    \item{Have a \textbf{Statement of Responsibilities} that clearly documents what they are personally responsible and accountable for.}
    \item{Take \textbf{reasonable steps} to prevent regulatory breaches in their area of responsibility — and to be held to account where they fail to do so.}
    \item{Comply with the \textbf{Conduct Rules} — the FCA's standards of behaviour for regulated individuals.}
\end{itemize}
\paragraph{The \textbf{"reasonable steps" standard} is practically significant — it creates a presumption that where a regulatory breach occurs in a senior manager's area of responsibility, they are accountable for it unless they can demonstrate that they took reasonable steps to prevent it. This reversal of the burden of proof — compared to the previous regime where regulators had to prove individual culpability — was designed to create a more powerful personal accountability incentive for senior leaders.}
\paragraph{For compliance professionals, the SM\&CR has several direct implications. The Chief Compliance Officer is typically an SMF holder — with personal accountability for the firm's compliance framework. The compliance function is responsible for supporting the firm's implementation of the regime — including maintaining accurate Statements of Responsibilities, supporting the certification process, and ensuring that the Conduct Rules are effectively communicated and embedded.}

\subsection{The Certification Regime}
\paragraph{The \textbf{Certification Regime} applies to a broader population of individuals — those performing \textbf{Certified Functions} whose roles mean that they could cause significant harm to the firm or its customers. Certified Function holders are not approved by regulators — the firm itself is responsible for assessing their fitness and propriety and issuing annual certifications confirming that they meet the required standard.}
\paragraph{The certification obligation creates a direct people risk assessment requirement at the individual level — firms must assess, at least annually, whether each certified function holder is fit and proper, taking into account:}
\begin{itemize}
    \bitem{Honesty, integrity, and reputation}{the individual's personal conduct, criminal record, regulatory history, and financial soundness.}
    \bitem{Competence and capability}{the individual's knowledge, skills, and experience relative to the requirements of their role.}
    \bitem{Financial soundness}{the individual's personal financial position — recognising that financial difficulties can create vulnerability to temptation or coercion.}
\end{itemize}
\paragraph{The certification process is a structured individual-level conduct and people risk assessment — and its quality determines whether the firm has genuine assurance that those performing significant roles meet the required standard.}

\subsection{The Conduct Rules}
\paragraph{The \textbf{Conduct Rules} are the FCA's standards of behaviour applicable to virtually all employees of regulated firms — not just Senior Managers and Certified Function holders. They establish five individual Conduct Rules:}
\begin{itemize}
    \bitem{Rule 1}{You must act with integrity.}
    \bitem{Rule 2}{You must act with due skill, care, and diligence.}
    \bitem{Rule 3}{You must be open and cooperative with the FCA, PRA, and other regulators.}
    \bitem{Rule 4}{You must pay due regard to the interests of customers and treat them fairly.}
    \bitem{Rule 5}{You must observe proper standards of market conduct.}
\end{itemize}
\paragraph{XXXX.}
\begin{itemize}
    \bitem{SM1}{You must take reasonable steps to ensure that the business of the firm for which you are responsible is organised so that it can be controlled effectively.}
    \bitem{SM2}{You must take reasonable steps to ensure that the business of the firm for which you are responsible complies with the relevant requirements and standards of the regulatory system.}
    \bitem{SM3}{You must take reasonable steps to ensure that any delegation of your responsibilities is to an appropriate person and that you oversee the discharge of the delegated responsibility effectively.}
    \bitem{SM4}{You must disclose appropriately any information of which the FCA or PRA would reasonably expect notice.}
\end{itemize}
\paragraph{The Conduct Rules are significant for people and conduct risk assessment because they establish clear, enforceable standards against which individual behaviour can be assessed — and because breaches of the Conduct Rules can be the basis for regulatory action against individuals.}

%---------------------------- SECTION ----------------------------
\section{Human Factors — The Science Behind People Risk}
\paragraph{To assess people and conduct risk effectively, it is helpful to understand the scientific foundation — the field of \textbf{human factors}, which studies how the characteristics of human beings interact with the systems, environments, and organisations in which they work to produce outcomes ranging from excellent to catastrophic.}
\paragraph{Human factors research has identified several consistent patterns in how human beings behave in complex organisational systems that are directly relevant to people risk assessment.}

\subsection{Error Types — Understanding How Mistakes Happen}
\paragraph{The most widely adopted framework for understanding human error in organisational contexts is \textbf{James Reason's taxonomy} — which distinguishes between different types of human failure at different levels of the performance system.}
\paragraph{\textbf{Slips and lapses} are failures of execution — where the intention is correct but the execution goes wrong. A slip is an action that was not intended — pressing the wrong key, misreading a number, transposing digits. A lapse is a failure of memory — forgetting a step in a procedure, losing track of where you are in a complex task. Slips and lapses are particularly common in tasks that are routine, automated, and performed under time pressure or distraction — which describes a significant proportion of the work done in most organisations.}
\paragraph{\textbf{Mistakes} are failures of planning or intention — where the action is executed as intended but the intention itself is wrong. Rule-based mistakes occur when an individual applies the wrong rule to a situation — either because they misclassify the situation, or because the rule they apply is itself flawed. Knowledge-based mistakes occur when an individual is operating in unfamiliar territory — beyond the scope of their established rules and procedures — and their reasoning leads them to an incorrect conclusion.}
\paragraph{\textbf{Violations} are intentional departures from required procedures or standards — where the individual knows what the rule is and chooses not to follow it. Violations range from routine violations — habitual shortcuts that have become normalised in a particular work environment — through situational violations — departures from procedure driven by specific pressures or circumstances — to exceptional violations — deliberate rule-breaking in the belief that it is necessary or justified. And at the far end, sabotage — intentional acts designed to cause harm.}
\paragraph{The distinction between violations and other error types is important for both risk assessment and control design. Slips, lapses, and mistakes can be addressed through improved process design, better system interfaces, more effective training, and enhanced monitoring. Violations require a different response — addressing the cultural, incentive, and accountability structures that make rule-breaking seem acceptable or attractive.}

\subsection{The Swiss Cheese Model — How Defences Fail}
\paragraph{\textbf{Reason's Swiss Cheese Model} of organisational accidents provides a powerful framework for understanding how serious harm occurs despite the existence of multiple defensive layers.}
\paragraph{In the model, each defensive layer — each control, each safeguard, each oversight mechanism — is represented as a slice of Swiss cheese. Each slice has holes — the weaknesses, gaps, and failures that mean no single defensive layer is perfect. In normal operation, the holes in different slices are not aligned — even if one defence fails, others catch the failure. A serious accident or harm occurs when, on a particular occasion, the holes in all the slices align — creating a trajectory through all the defensive layers from hazard to harm.}
\paragraph{The Swiss Cheese model has several important implications for people risk assessment:}
\paragraph{\textbf{Single points of failure are dangerous} — a defensive system that relies on a single human layer — a single reviewer, a single approval, a single check — is vulnerable to the failure of that single layer in a way that a multi-layered system is not.}
\paragraph{\textbf{Latent failures are as important as active failures} — Reason distinguishes between \textbf{active failures} — the immediate actions or errors that trigger a harmful event — and \textbf{latent failures} — the underlying weaknesses in the system that created the conditions for the active failure to have serious consequences. Most people risk assessments focus too heavily on active failures and pay insufficient attention to the latent conditions — the inadequate training, the excessive workload, the missing procedure, the cultural norm that makes shortcuts acceptable — that make active failures dangerous.}
\paragraph{\textbf{The system shapes individual behaviour} — the model's most important insight for compliance professionals may be this: the behaviour of individuals is substantially shaped by the systems, processes, and cultures within which they work. Addressing people risk by focusing only on the individuals — improving their training, increasing their supervision, removing the worst performers — without addressing the systemic conditions that create the risk is unlikely to produce lasting improvement.}

\subsection{Cognitive Biases and Decision-Making Failures}
\paragraph{Beyond the error taxonomy, human factors research has documented a wide range of cognitive biases — systematic patterns of deviation from rational judgement — that affect decision-making in organisational contexts. Several are particularly relevant to conduct and people risk assessment:}
\paragraph{\textbf{Confirmation bias} — the tendency to seek out and give greater weight to information that confirms existing beliefs whilst discounting contradictory evidence. In a compliance context, confirmation bias can lead managers to overlook warning signs that contradict a positive assessment of an individual or a team, and can cause compliance professionals to look for evidence that a practice is compliant rather than genuinely testing whether it is.}
\paragraph{\textbf{Authority bias} — the tendency to defer to the judgements of authority figures, even when those judgements conflict with one's own assessment or with the available evidence. In organisational hierarchies, authority bias can suppress challenge and dissent — particularly when senior individuals express strong views, as noted in Chapter~\ref{chap:Communicating Risk Across the Organisation}'s discussion of groupthink.}
\paragraph{\textbf{Normalisation of deviance} — the gradual process by which departures from required standards become normalised through repeated occurrence without apparent adverse consequence. The concept — developed by sociologist Diane Vaughan in her analysis of the Challenger space shuttle disaster — describes how organisations can drift, incrementally and almost invisibly, from compliant to non-compliant behaviour as minor violations accumulate and become the new baseline.}
\paragraph{\textbf{Overconfidence} — the well-documented tendency to overestimate one's own competence, the reliability of one's judgements, and the adequacy of one's preparations. Overconfidence is particularly dangerous in senior individuals — who may have had their confidence reinforced by a track record of success — and in novel or rapidly changing environments where past experience provides a less reliable guide to future performance.}
\paragraph{\textbf{Moral disengagement} — the psychological processes through which individuals disengage their normal ethical standards from their behaviour — justifying harmful actions through rationalisation, displacement of responsibility, dehumanisation of those affected, or diffusion of responsibility in group settings. Understanding moral disengagement is essential for assessing conduct risk — particularly in cultures where boundary-pushing is normalised and where group dynamics suppress individual ethical judgement.}

%---------------------------- SECTION ----------------------------
\section{Assessing People Risk — A Structured Framework}
\paragraph{With the conceptual foundations in place, what does a practical people risk assessment look like? The framework applies the core five-step process to the specific characteristics of people risk.}

\subsection{Identifying People Risk Sources}
\paragraph{People risk sources can be identified across several dimensions:}
\paragraph{\textbf{Role-level analysis} — examining each significant role in the organisation and identifying the people risks associated with it. Which roles involve consequential decisions with limited oversight? Which roles involve handling of sensitive data or assets that create fraud or theft risk? Which roles involve customer contact that creates conduct risk? Which roles are dependent on individual expertise that is not documented or transferable?}
\paragraph{\textbf{Process-level analysis} — examining key processes for the human error and conduct risks embedded in their design. Where does the process rely on individual judgement without adequate guidance? Where are manual steps that could be prone to error? Where does the process create time pressure that increases error probability? Where does the design of the process create incentives for cutting corners?}
\paragraph{\textbf{Cultural analysis} — examining the organisation's culture for the conditions that enable or encourage poor conduct. What behaviours are actually rewarded — as opposed to those that are formally valued? Are there subcultures within specific teams or functions where norms diverge from the organisation's stated standards? Are there known practices that are inconsistent with policy but have been tolerated because they are commercially productive?}
\paragraph{\textbf{Incentive analysis} — examining the incentive structures that shape individual behaviour. Where do financial incentives — bonuses, commissions, targets — create pressure that could drive conduct risk? Where do career incentives — promotion, recognition, status — reward outcomes without adequate attention to how they are achieved? Where do negative incentives — fear of failure, job insecurity, excessive performance pressure — create conditions for shortcuts, concealment, or corner-cutting?}
\paragraph{\textbf{Workforce analysis} — examining the composition and characteristics of the workforce for people risk indicators. Where are key person dependencies most acute? Where are skill gaps most significant? Where is turnover highest — and what does that signal about the culture and conditions in those areas? Where are workloads most excessive — creating the time pressure and fatigue that increase error probability?}

\subsection{Determining Who Is Affected}
\paragraph{The affected parties analysis for people risk has several dimensions:}
\paragraph{\textbf{Customers} — the most significant external affected party for conduct risk. Poor individual conduct — misselling, poor advice, unfair treatment, deliberate deception — directly harms customers and creates regulatory and legal exposure.}
\paragraph{\textbf{Colleagues and the organisation} — internal misconduct — fraud, harassment, bullying, unauthorised disclosure — harms colleagues and the organisation directly. Inadequate individual competence creates risk for the colleagues who depend on that person's work being done correctly.}
\paragraph{\textbf{Markets and counterparties} — market conduct risk — insider trading, benchmark manipulation, market abuse — creates harm to market participants and undermines market integrity.}
\paragraph{\textbf{The organisation's regulatory standing} — conduct failures create direct regulatory risk — for the organisation and, under the SM\&CR, for the individuals involved.}

\subsection{Evaluating People Risk}
\paragraph{People risk evaluation involves assessing both the likelihood of the risk materialising — the probability that the identified people risk will actually result in harmful conduct — and the potential consequence — the severity of the harm if it does.}
\paragraph{Specific evaluation considerations for people risk include:}
\paragraph{\textbf{Control environment quality} — the strength of the controls designed to prevent, detect, and respond to people risk events. The more robust the control environment — the more effective the supervision, the more rigorous the performance management, the more meaningful the oversight — the lower the residual likelihood that harmful conduct will occur undetected.}
\paragraph{\textbf{Cultural indicators} — the cultural signals that increase or decrease the probability of conduct risk materialising. High staff turnover in a specific area, patterns of overriding controls, excessive pressure on performance metrics, and absence of psychological safety are all cultural indicators that increase conduct risk likelihood.}
\paragraph{\textbf{Historical experience} — the organisation's own incident and near-miss history in the relevant risk area. Past conduct failures — even those that did not result in significant harm — are indicators of underlying risk that is likely to persist if the underlying conditions are not addressed.}
\paragraph{\textbf{Individual risk indicators} — for certified function assessments and individual-level people risk assessment, specific indicators of individual risk include regulatory history, financial soundness concerns, unexplained lifestyle changes, unusual patterns of behaviour, and engagement with high-risk activities outside the organisation.}

\subsection{Controls for People Risk}
\paragraph{The controls available for managing people risk span the full range of the control hierarchy — from the elimination of risk through process redesign, through engineering controls embedded in systems, to the administrative controls of policy, training, and supervision.}
\paragraph{\textbf{Recruitment and selection} — the quality of the people who join the organisation is the first line of defence against people risk. Effective screening — including criminal record checks, financial soundness checks, regulatory history checks, and reference checks that genuinely probe past conduct — reduces the probability of appointing individuals who carry significant conduct risk.}
\paragraph{\textbf{Role design and segregation of duties} — designing roles and processes to reduce the opportunity for harmful conduct through appropriate segregation of duties, mandatory oversight of consequential decisions, and limits on individual authority that are calibrated to the risk of the decisions involved.}
\paragraph{\textbf{Training and competence frameworks} — ensuring that individuals have the knowledge and skills to perform their roles to the required standard, and that they understand the conduct standards they are expected to meet. Training needs to be genuine — genuinely informative, genuinely relevant, and genuinely assessed — not a completion exercise.}
\paragraph{\textbf{Supervision and oversight} — active, genuine supervision of conduct — particularly in high-risk roles and high-risk areas — is one of the most effective conduct risk controls. This means supervisors who actually observe the work being done, who review outputs for quality and compliance, and who are alert to the early warning signs of conduct risk. In financial services, supervision of customer-facing staff is a specific regulatory expectation.}
\paragraph{\textbf{Performance management} — performance management frameworks that genuinely assess how results are achieved, not just whether they are achieved. This means incorporating conduct, compliance, and risk management into performance assessments — rewarding good conduct and imposing meaningful consequences for poor conduct}
\paragraph{\textbf{Whistleblowing arrangements} — effective, accessible, and genuinely protected whistleblowing channels that enable individuals to raise concerns about the conduct of others. The quality of the organisation's whistleblowing arrangements is both a conduct risk control — enabling early detection of misconduct — and a cultural indicator — reflecting the degree to which the organisation genuinely values the surfacing of concerns.}
\paragraph{\textbf{Monitoring and surveillance} — systematic monitoring of high-risk conduct indicators — communications surveillance, transaction monitoring, management information analysis — that enables the detection of conduct issues before they escalate. For financial services firms, the surveillance of trading activity and communications is a specific regulatory expectation.}
\paragraph{\textbf{Disciplinary and accountability frameworks} — clear, consistently applied consequences for conduct failures. Nothing undermines a conduct risk framework more completely than a culture in which known conduct failures are tolerated because the individuals involved are commercially valuable, or where accountability is applied selectively based on seniority or performance. Consistency of consequence — across all levels of the organisation — is a primary determinant of cultural credibility.}

%---------------------------- SECTION ----------------------------
\section{Specific People Risk Areas — A Closer Look}

\subsection{Key Person Risk}
\paragraph{\textbf{Key person risk} — the operational and knowledge risk created by over-dependence on specific individuals — is a pervasive and frequently underestimated category of people risk. In most organisations, there are individuals whose absence — through resignation, illness, incapacity, or death — would create significant operational disruption, knowledge loss, or relationship damage.}
\paragraph{Key person risk assessment involves:}
\begin{itemize}
    \item{Identifying roles and individuals where the risk is most acute — where the knowledge, relationships, or skills involved are not documented, not shared, and not easily replaced.}
    \item{Assessing the consequences of the key person's sudden unavailability — the operational, regulatory, and relationship impacts.}
    \item{Identifying mitigation measures — succession planning, knowledge documentation, cross-training, and relationship broadening — and assessing their adequacy and implementation status.}
\end{itemize}
\paragraph{For regulated firms, the SM\&CR's requirement for accurate Statements of Responsibilities helps to mitigate key person risk by ensuring that responsibilities are documented and assigned — but it does not substitute for genuine succession planning and knowledge transfer.}

\subsection{Fraud and Dishonesty Risk}
\paragraph{\textbf{Internal fraud} — the deliberate theft, misappropriation, or misuse of assets, information, or authority by employees — is a significant and persistent people risk category. The \textbf{Association of Certified Fraud Examiners (ACFE)} estimates that organisations typically lose around five percent of annual revenues to occupational fraud — a figure that, whilst difficult to validate precisely, is consistently supported by loss data from those who have experienced it.}
\paragraph{The \textbf{Fraud Triangle} — developed by criminologist Donald Cressey — provides a widely adopted framework for understanding the conditions under which individuals are likely to commit fraud. The three elements are:}
\paragraph{\textbf{Pressure} — a financial or personal pressure that motivates the individual to commit fraud. This might be personal financial difficulty, addiction, lifestyle aspirations beyond legitimate income, or pressure to meet performance targets that are unachievable through legitimate means.}
\paragraph{\textbf{Opportunity} — access to assets, information, or systems combined with the ability to conceal the fraud. Weak controls, inadequate segregation of duties, excessive trust, and poor oversight all create opportunity.}
\paragraph{\textbf{Rationalisation} — the psychological justification that enables the individual to reframe dishonest behaviour as acceptable — \textit{"the organisation owes me", "I'll pay it back", "everyone does it", "they'll never notice"}.}
\paragraph{The Fraud Triangle is a practical assessment tool — because it suggests that fraud risk can be reduced by addressing any of the three elements. Reducing pressure — through employee financial wellness programmes and by not setting unachievable targets — reduces motivation. Strengthening controls and reducing opportunity — through segregation of duties, mandatory holidays for key roles, and rotation of high-risk positions — limits the practical ability to commit fraud. And building a culture of integrity — where rationalisation is harder because the organisation's values are genuinely lived — reduces the psychological availability of the justifications that enable dishonest behaviour.}

\subsection{Harassment, Bullying, and Workplace Culture Risk}
\paragraph{The risk of harassment, bullying, and discriminatory conduct within the workplace creates harm to employees, creates legal and regulatory exposure for the organisation, and is a direct indicator of cultural conditions that are likely to create other conduct risks. People who are treated badly are less likely to raise concerns, less likely to report misconduct, and less likely to maintain the ethical standards that an effective conduct risk framework requires.}
\paragraph{Assessing workplace culture risk involves:}
\begin{itemize}
    \item{Monitoring employee experience data — engagement surveys, exit interview data, grievance and disciplinary trends — for indicators of cultural problems.}
    \item{Understanding the pattern of formal complaints and how they are handled — whether they are taken seriously, whether they result in appropriate outcomes, and whether complainants are protected from retaliation.}
    \item{Being alert to the informal cultural signals — the team where nobody ever challenges the manager, the office where jokes about customers are normalised, the leader whose track record of delivering results has made their interpersonal conduct exempt from scrutiny.}
\end{itemize}

%---------------------------- SECTION ----------------------------
\section{The SM\&CR and Individual Conduct Assessment in Practice}
\paragraph{For compliance professionals in regulated financial services firms, the SM\&CR creates a direct and ongoing obligation to assess individual conduct — through the certification process, through Conduct Rules training and monitoring, and through the assessment and management of fitness and propriety concerns.}
\paragraph{In practice, effective SM\&CR compliance involves:.}
\paragraph{\textbf{Maintaining a complete and accurate population of Senior Managers and Certified Function holders} — understanding who holds SMFs and CFs at any given time, ensuring that the regulatory position accurately reflects the actual population, and managing changes promptly and accurately.}
\paragraph{\textbf{Conducting genuine, evidence-based fitness and propriety assessments} — not simply relying on self-certification or on the absence of obvious red flags, but actively gathering and evaluating evidence of fitness and propriety at the point of appointment and at annual renewal.}
\paragraph{\textbf{Maintaining Statements of Responsibilities that genuinely reflect the distribution of responsibilities} — and reviewing them promptly when responsibilities change, to ensure that the accountability framework remains accurate and credible.}
\paragraph{\textbf{Delivering Conduct Rules training that genuinely develops understanding} — not simply completing a training module that confirms awareness of the rules' existence, but engaging people with what the rules mean in practice and how they apply to the specific situations they encounter in their roles.}
\paragraph{\textbf{Maintaining a process for assessing, recording, and responding to potential Conduct Rules breaches} — including a clear escalation path for concerns about individual conduct, a process for investigating those concerns proportionately and fairly, and a mechanism for regulatory notification where a breach is identified.}

%---------------------------- SECTION ----------------------------
\newpage
\section{Chapter Summary}
In this chapter we learned about the following key points:

\begin{table}[H]
  \centering
  \renewcommand{\arraystretch}{1.5}
  \begin{tabularx}{\textwidth}{ | >{\raggedright\arraybackslash}p{0.3\textwidth} 
								| >{\raggedright\arraybackslash}p{0.35\textwidth} 
                                | >{\raggedright\arraybackslash}X | }
    \hline
    \textbf{Risk Type} & \textbf{Key Sources} & \textbf{Primary Controls} \\
    \hline
    \textbf{Human error} & Task design; time pressure; fatigue; distraction; inadequate training & Process redesign; system automation; error-catching controls; training\\
    \hline
    \textbf{Inadequate competence} & Recruitment quality; training adequacy; role design & Competency frameworks; training; supervision; certification\\
    \hline
    \textbf{Conduct risk — customer} & Incentive misalignment; culture; inadequate supervision & Conduct framework; supervision; outcome monitoring; SM\&CR\\
    \hline
    \textbf{Conduct risk — market} & Opportunity; culture; inadequate surveillance & Surveillance; segregation; conflict management; SM\&CR\\
    \hline
    \textbf{Internal fraud} & Fraud Triangle — pressure, opportunity, rationalisation & Segregation of duties; rotation; anti-fraud controls; culture\\
    \hline
    \textbf{Key person risk} & Knowledge concentration; inadequate succession & Documentation; cross-training; succession planning\\
    \hline
    \textbf{Workplace culture risk} & Leadership behaviour; incentives; accountability gaps & Culture assessment; leadership accountability; grievance processes\\
    \hline
  \end{tabularx}
  \caption{Types of Risk In Regulated Industry}
\end{table}

\begin{table}[H]
  \centering
  \renewcommand{\arraystretch}{1.5}
  \begin{tabularx}{\textwidth}{ | >{\raggedright\arraybackslash}p{0.35\textwidth} 
								| >{\raggedright\arraybackslash}p{0.35\textwidth} 
                                | >{\raggedright\arraybackslash}X | }
    \hline
    \textbf{Element} & \textbf{Risk Indicators} & \textbf{Mitigation Approach} \\
    \hline
    \textbf{Pressure} & Financial difficulty; unachievable targets; personal stress & Employee support; realistic targets; financial wellness\\
    \hline
    \textbf{Opportunity} & Weak controls; poor segregation; inadequate oversight & Segregation; rotation; monitoring; access controls\\
    \hline
    \textbf{Rationalisation} & Cultural acceptance of rule-bending; perceived unfairness & Culture; values; consistent accountability; ethical leadership\\
    \hline
  \end{tabularx}
  \caption{The Fraud Triangle}
\end{table}

\begin{table}[H]
  \centering
  \renewcommand{\arraystretch}{1.5}
  \begin{tabularx}{\textwidth}{ | >{\raggedright\arraybackslash}p{0.35\textwidth} 
								| >{\raggedright\arraybackslash}p{0.35\textwidth} 
                                | >{\raggedright\arraybackslash}X | }
    \hline
    \textbf{Component} & \textbf{Key Obligation} & \textbf{Compliance Responsibility} \\
    \hline
    \textbf{Senior Managers Regime} & Approved individuals; Statements of Responsibilities; reasonable steps & Maintain accurate SMF population; SoR quality; accountability framework\\
    \hline
    \textbf{Certification Regime} & Annual fitness and propriety assessment & Genuine evidence-based certification process\\
    \hline
    \textbf{Conduct Rules} & Training; monitoring; breach reporting & Effective training; breach identification; regulatory notification\\
    \hline
  \end{tabularx}
  \caption{Senior Management And Certification Regime}
\end{table}

\paragraph{Key takeaways:}
\begin{itemize}
  \item{People are the primary source of most organisational risk — addressing people and conduct risk is the most fundamental investment in overall risk management effectiveness.}
  \item{Human error is systemic as well as individual — the conditions that produce error are as important as the errors themselves, and the system must be assessed and redesigned as well as the individuals trained.}
  \item{Conduct risk arises from the interaction of individual behaviour, cultural conditions, and incentive structures — addressing it requires all three to be assessed and managed.}
  \item{The Fraud Triangle provides a practical framework for understanding and mitigating internal fraud risk — addressing pressure, opportunity, and rationalisation together.}
  \item{The SM\&CR creates a direct, ongoing individual-level conduct risk assessment obligation for regulated financial services firms — certification is a risk assessment exercise, not an administrative one.}
  \item{Normalisation of deviance — the incremental drift from compliant to non-compliant behaviour — is one of the most dangerous and most insidious people risk dynamics, and detecting it requires cultural as well as process monitoring.}
  \item{Consistent accountability — consequences applied fairly and at all levels — is the single most important determinant of a culture in which conduct risk is genuinely managed rather than merely addressed on paper.}
\end{itemize}

\barquote{In Chapter~\ref{chap:Dynamic and Continuous Risk Assessment}, the final chapter of Part VI, we turn to the most strategically significant category of risk — strategic and emerging risk assessment. We explore how organisations assess risks that operate at the level of strategy and business model viability, how to identify and assess risks that do not yet exist, and how the compliance and legal function contributes to the strategic risk conversations that determine the organisation's long-term resilience and relevance.}